% File:
%   - malbolge-specific-optimization-mathematics.tex
% Path:
%   - math/algorithms/malbolge-specific-optimization-mathematics.tex
%
% Copyright:
%   - Copyright (c) 2026 Alberto Villa Osorno.
% SPDX-License-Identifier:
%   - MIT
% Confidential:
%   - false
% License-File:
%   - LICENSE
% Path-Rule:
%   - All paths in this header are repository-root relative.
%
% Boundary-Contract:
% - Owns:
%   - The repository behavior implemented by this source file.
% - Must-Not:
%   - Bypass the contracts or authority boundaries of its owning package.
% - Allows:
%   - Inputs: values admitted by the file's public or internal interface.
%   - Outputs: deterministic values or effects declared by that interface.
%   - Side effects: only those explicitly owned by the implementation.
% - Split-When:
%   - Split when one responsibility gains an independent lifecycle.
% - Merge-When:
%   - Merge when another file owns the exact same responsibility.
% - Summary:
%   - Malbolge specific optimization mathematics module.
% - Description:
%   - Implements the responsibility summarized by this module.
% - Usage:
%   - Used through the owning package, executable, or document boundary.
% - Defaults:
%   - Invalid inputs or broken invariants fail closed.
%
% Related documents:
% - None.
%
% Large file:
%   - false
%

%! Malbolge specific optimization mathematics module.

\documentclass{article}
\usepackage{amsmath,amssymb}
% Copyright (c) 2026 Alberto Villa Osorno.
% SPDX-License-Identifier: MIT
% Shared notation only; this file is included by buildable math documents.

\newcommand{\TritSet}{\{0,1,2\}}
\newcommand{\WordModulus}[1]{3^{#1}}
\newcommand{\WordDomain}[1]{\mathcal{W}_{#1}}
\newcommand{\MemoryState}[1]{\mathcal{M}_{#1}}
\newcommand{\CrazyDigit}{\gamma}
\newcommand{\CrazyN}[3]{\Gamma_{#1}\!\left(#2,#3\right)}
\newcommand{\RotateN}[2]{\rho_{#1}\!\left(#2\right)}
\newcommand{\DecodeOp}{\delta}
\newcommand{\EncryptOp}{\varepsilon}
\newcommand{\LoadOp}{\mathsf{Load}}
\newcommand{\LowerOp}{\mathsf{Lower}}
\newcommand{\VerifyOp}{\mathsf{Verify}}
\newcommand{\ObservedEquivalent}[1]{\equiv_{\mathrm{obs},#1}}
\newcommand{\CostVector}[1]{\mathbf{K}\!\left(#1\right)}
\newcommand{\Transition}[1]{\xrightarrow{#1}}

\title{Malbolge-Specific Optimization Mathematics}
\author{Malbolge Project}
\date{2026-07-27}
\begin{document}
\maketitle

\section{Scope}

This document records algebraic reductions that are already consumed by the
Malbolge runtime. A reduction is admitted only on its stated finite/profile
domain and only while its correspondence evidence remains valid.

\section{Classic crazy chunk factorization}

Let $B=3^5=243$. For classic ten-trit words write
\[
  d=d_0+B d_1,\qquad a=a_0+B a_1,
\]
with $d_0,d_1,a_0,a_1\in\{0,\ldots,B-1\}$. Because crazy is digitwise,
\begin{equation}
  \CrazyN{10}{d}{a}
    = \CrazyN{5}{d_0}{a_0}
      + B\CrazyN{5}{d_1}{a_1}.
  \label{eq:optimization-crazy-chunk}
\end{equation}
Therefore one exact $243\times243=59049$ five-trit table is sufficient for both
halves of every classic crazy operation; a $59049^2$ word-pair table is neither
required nor justified.

\section{Profile-width crazy chunk factorization}

For profile width $N$, let chunk width $k=5$, chunk index $j\geq0$, and
\[
  r_j=\min(k,N-jk).
\]
For every admitted chunk with $r_j>0$, define
\[
  d_j=\left\lfloor\frac{d}{3^{jk}}\right\rfloor\bmod3^{r_j},
  \qquad
  a_j=\left\lfloor\frac{a}{3^{jk}}\right\rfloor\bmod3^{r_j}.
\]
Digit independence gives
\begin{equation}
  \CrazyN{N}{d}{a}
    = \sum_{j=0}^{\lceil N/k\rceil-1}
      3^{jk}\left(\CrazyN{5}{d_j}{a_j}\bmod3^{r_j}\right).
  \label{eq:optimization-crazy-profile-chunks}
\end{equation}
The current 14-trit runtime therefore uses chunks $5+5+4$ while reusing the same
five-trit table. The final four-trit result is reduced modulo $3^4$; no phantom
fifteenth trit is admitted.

\section{Decode phase reduction}

For graphical cell $m$, decode depends on the code pointer only through its
residue modulo 94:
\begin{equation}
  \DecodeOp(m,c)=\DecodeOp(m,c\bmod94).
  \label{eq:optimization-decode-phase}
\end{equation}
The runtime may therefore precompute one 94-value code phase and a
$94\times94$ graphical-cell/phase table without changing decode semantics.

\section{Classic rotate lookup}

The classic word domain is finite, so define a table $R$ by
\begin{equation}
  R[x]=\RotateN{10}{x}
  \qquad\text{for every }x\in\WordDomain{10}.
  \label{eq:optimization-rotate-table}
\end{equation}
Replacing scalar rotation by $R[x]$ is an exact finite-domain substitution, not
an approximation or heuristic.

\section{Verified candidate boundary}

Let $C$ be a candidate set, $V(c,p)\in\{0,1\}$ the independent verifier under
profile $p$, and $\CostVector{c}$ the measured cost vector. Search may rank or
generate candidates arbitrarily, but only
\[
  C_V=\{c\in C\mid V(c,p)=1\}
\]
is eligible for cost comparison. A speedup cannot repair failed semantic
correspondence.

\section{Measured implementation result}

The versioned CPU VM benchmark evidence at
\texttt{benchmarks/interpreter/evidence/2026-07-26-windows-x86\_64/} contains 15 raw
samples per scalar/table implementation and identical checksums. On that
identified host, the classic crazy factorization reduced the median from
77,456,700 ns to 7,423,600 ns (10.43x), while rotate lookup reduced the median
from 15,260,300 ns to 10,141,700 ns (1.50x). These measurements justify the
current CPU implementation on that recorded workload; they are not universal
hardware performance claims.

\end{document}
